\documentclass[12pt, a4paper]{article}
\usepackage[UTF8]{ctex}
\usepackage{amsmath, amssymb, amsthm}
\usepackage{geometry}
\usepackage{bm}

\geometry{left=2.5cm, right=2.5cm, top=2.5cm, bottom=2.5cm}

\newcommand{\RR}{\mathbb{R}}
\newcommand{\e}{\mathrm{e}}
\newcommand{\dif}{\mathrm{d}}

\begin{document}
\section*{13.3}
\textbf{13.}证明不等式
\[2\pi(\sqrt{17} - 4) \leq \iint\limits_{x^2 + y^2 \leq 1} \frac{\dif x \dif y}{\sqrt{16 + \sin^2 x + \sin^2 y}} \leq \frac{\pi}{4}.\]

证明：
\[\iint\limits_{x^2 + y^2 \leq 1} \frac{\dif x \dif y}{\sqrt{16 + \sin^2 x + \sin^2 y}} \leq \iint\limits_{x^2 + y^2 \leq 1} \frac{\dif x \dif y}{4} = \frac{\pi}{4}.\]
\[\iint\limits_{x^2 + y^2 \leq 1} \frac{\dif x \dif y}{\sqrt{16 + \sin^2 x + \sin^2 y}} \geq \iint\limits_{x^2 + y^2 \leq 1} \frac{\dif x \dif y}{\sqrt{16 + x^2 + y^2}} = \int_{0}^{2\pi} \dif \theta \int_0^1 \frac{r\dif r}{\sqrt{16 + r^2}} = 2\pi(\sqrt{17} - 4).\]

\textbf{14.}设一元函数\(f(u)\)在\([-1, 1]\)上连续，证明：
\[\iint\limits_{|x| + |y| \leq 1} f(x + y) \dif x \dif y = \int_{-1}^{1} f(u) \dif u.\]

证明：令\(u = x + y, v = x - y\).
\[J = \left|\frac{\partial(x, y)}{\partial(u, v)}\right| = -\frac{1}{2}.\]
\[\iint\limits_{|x| + |y| \leq 1} f(x + y) \dif x \dif y = \frac{1}{2} \iint\limits_{\max\{|u|, |v|\} \leq 1} f(u) \dif u \dif v = \int_{-1}^{1} f(u) \dif u.\]

\textbf{16.}计算下列\(n\)重积分：

(1)\(\int\limits_\Omega \sqrt{x_1 + x_2 + \cdots + x_n} \dif x_1 \dif x_2 \cdots \dif x_n\)，其中\\
\(\Omega = \{(x_1, x_2, \cdots, x_n)|x_1 + x_2 + \cdots + x_n \leq 1, x_i \geq 0, i = 1, 2, \cdots, n\}\).

解：令\(\begin{cases}
y_1 = x_1 + x_2 + \cdots + x_n, \\
y_2 = x_2  + \cdots + x_n, \\
\cdots \\
y_n = x_n,
\end{cases}\)，则
\[\int\limits_\Omega \sqrt{x_1 + x_2 + \cdots + x_n} \dif x_1 \dif x_2 \cdots \dif x_n = \int_0^1 \sqrt{y_1} \dif y_1 \int_{0}^{y_1} \dif y_2 \cdots \int_{0}^{y_{n - 1}} \dif y_n = \frac{2}{(n - 1)!(2n + 1)}.\]

(2)\(\int\limits_\Omega (x_1^2 + x_2^2 + \cdots + x_n^2) \dif x_1 \dif x_2 \cdots \dif x_n\)，\\
其中\(\Omega\)为\(n\)维球体\(\{(x_1, x_2, \cdots, x_n)|x_1^2 + x_2^2 + \cdots + x_n^2 \leq 1\}\).

解：
\[\int\limits_\Omega (x_1^2 + x_2^2 + \cdots + x_n^2) \dif x_1 \dif x_2 \cdots \dif x_n = \int_0^1 r^2 \cdot r^{n - 1} \dif r \int\limits_{\Omega_{n - 1}} \dif \Omega_{n - 1} = \frac{2\pi^{\frac{n}{2}}}{(n + 2)\left(\dfrac{n}{2}\right)!}。\]

\section*{13.4}
\textbf{1.}讨论下列反常重积分的敛散性：

(1)\(\iint\limits_{\RR^2} \dfrac{\dif x \dif y}{(1 + |x|^p)(1 + |y|^q)}\).

解：
\begin{align*}
&\iint\limits_{\RR^2} \frac{\dif x \dif y}{(1 + |x|^p)(1 + |y|^q)} \\
= &\int_{-\infty}^{+\infty} \frac{\dif x}{(1 + |x|^p)} \int_{-\infty}^{+\infty} \frac{\dif y}{(1 + |y|^q)} \\
~ &\int_{-\infty}^{+\infty} \frac{\dif x}{|x|^p} \int_{-\infty}^{+\infty} \frac{\dif y}{|y|^p}
\end{align*}
故\(p > 1\)且\(q > 1\)时收敛.

(3)\(\iint\limits_{x^2 + y^2 \leq 1} \dfrac{\varphi(x, y)}{(1 - x^2 - y^2)^p} \dif x \dif y\)，其中\(\varphi(x, y)\)满足\\
\(0 < m \leq |\varphi(x, y)| \leq M\)（\(m, M\)为常数）.

解：令\(x = r\cos\theta, y = r\sin\theta\)，
\begin{align*}
&\iint\limits_{x^2 + y^2 \leq 1} \dfrac{\varphi(x, y)}{(1 - x^2 - y^2)^p} \dif x \dif y \\
= &\int_{0}^{2\pi} \dif \theta \int_0^1 \frac{r\varphi(r\cos\theta, r\sin\theta)}{(1 - r^2)^p}
\end{align*}
\[m \int_{0}^{2\pi} \dif \theta \int_0^1 \frac{r}{(1 - r^2)^p} \leq I \leq M \int_{0}^{2\pi} \dif \theta \int_0^1 \frac{r}{(1 - r^2)^p}\]
令\(u = 1 - r^2 \in [0, 1]\)，
\[\int_{0}^{2\pi} \dif \theta \int_0^1 \frac{r}{(1 - r^2)^p} = \frac{1}{2} \int_0^1 \frac{1}{u^p} \dif u.\]
当且仅当\(p < 1\)时收敛.

(4)\(\iint\limits_{[0, a] \times [0, a]} \dfrac{\dif x \dif y}{|x - y|^p}\).

解：令\(u = x - y, v = x + y, J = \dfrac{1}{2}\).
\[\iint\limits_{[0, a] \times [0, a]} \frac{\dif x \dif y}{|x - y|^p} = \frac{1}{2} \iint\limits_D \frac{\dif u \dif v}{|u|^p} = 2a \int_0^a u^{-p} \dif u - 2 \int_0^a u^{1 - p} \dif u.\]
当且仅当\(p < 1\)时收敛.

(5)\(\iiint\limits_{x^2 + y^2 + z^2 \leq 1} \dfrac{\dif x \dif y \dif z}{(x^2 + y^2 + z^2)^p}\).

解：令\(x = r\cos\theta, y = r\sin\theta\cos\varphi, z = r\sin\theta\sin\varphi\).
\[\iiint\limits_{x^2 + y^2 + z^2 \leq 1} \frac{\dif x \dif y \dif z}{(x^2 + y^2 + z^2)^p} = \int_{0}^{2\pi} \dif \varphi \int_0^\pi \sin\theta \dif \theta \int_0^1 r^{2 - 2p} \dif r = 4\pi \int_0^1 r^{2 - 2p} \dif r.\]
当且仅当\(p < \dfrac{3}{2}\)时收敛.

\textbf{2.}计算下列反常重积分：

(1)\(\iint\limits_D \dfrac{\dif x \dif y}{x^py^q}\)，其中\(D = \{(x, y)|xy \geq 1, x \geq 1\}\)，且\(p > q > 1\).

\begin{align*}
&\iint\limits_D \dfrac{\dif x \dif y}{x^py^q} \\
= &\int_{1}^{+\infty} \frac{\dif x}{x^p} \int_{\frac{1}{x}}^{+\infty} \frac{\dif y}{y^q} \\
= &\frac{1}{q - 1} \int_{1}^{+\infty} x^{-(p - q + 1)} \dif x \\
= &\frac{1}{(q - 1)(p - q)}.
\end{align*}

(3)\(\iiint\limits_{\RR^3} \e^{-(x^2 + y^2 + z^2)} \dif x \dif y \dif z\).

\begin{align*}
&\iiint\limits_{\RR^3} \e^{-(x^2 + y^2 + z^2)} \\
= &\int_{0}^{2\pi} \dif \varphi \int_0^\pi \sin\theta \dif \theta \int_{0}^{+\infty} e^{-r^2} \dif r \\
= &\pi^{\frac{3}{2}}.
\end{align*}

\textbf{4.}判断反常重积分
\[I = \iint\limits_{\RR^2} \frac{\dif x \dif y}{(1 + x^2)(1 + y^2)}\]
是否收敛.如果收敛，求其值.

解：
\[I = \int_{-\infty}^{+\infty} \frac{\dif x}{(1 + x^2)} \int_{-\infty}^{+\infty} \frac{\dif y}{(1 + y^2)} = \pi^2.\]

\textbf{6.}设函数\(f(x)\)在\([0, a]\)上连续，证明：
\[\iint\limits_{0 \leq y \leq x \leq a} \frac{f(y)}{\sqrt{(a - x)(x - y)}} \dif x \dif y = \pi \int_0^a f(x) \dif x.\]

解：
\begin{align*}
&\iint\limits_{0 \leq y \leq x \leq a} \frac{f(y)}{\sqrt{(a - x)(x - y)}} \dif x \dif y \\
= &\int_0^a f(y) \dif y \int_y^a \frac{1}{\sqrt{(a - x)(x - y)}} \dif x \\
= &\int_0^a f(y) \dif y \int_y^a \frac{1}{\sqrt{[a - y - (x - y)](x - y)}} \dif (x - y) \\
= &\int_0^a f(y) \dif y \int_{0}^{a - y} \frac{a - y}{\sqrt{[(a - y) - t]t}} \dif \frac{t}{a - y} \\
= &\int_0^a f(y) \dif y \int_0^1 \frac{1}{\sqrt{(1 - u)u}} \dif u \\
= &\int_0^a f(y) \dif y \int_{0}^{\frac{\pi}{2}} \frac{1}{\sqrt{(1 - \sin^2\theta)\sin^2\theta}} \dif (\sin^2\theta) \\
= &\pi\int_0^a f(y) \dif y.
\end{align*}

\textbf{7.}计算积分\(\int\limits_{\RR^n} \e^{-(x_1^2 + x_2^2 + \cdots + x_n^2)} \dif x_1 \dif x_2 \cdots \dif x_n\).

解：
\begin{align*}
&\int\limits_{\RR^2} \e^{-(x_1^2 + x_2^2 + \cdots + x_n^2)} \dif x_1 \dif x_2 \cdots \dif x_n \\
= &\left(\int_{-\infty}^{+\infty} \e^{-x^2} \dif x\right)^n \\
= &\pi^{\frac{n}{2}}.
\end{align*}

\end{document}